\chapter{Abstract}

\section{Abstract}

Machine learning (ML) models deployed in high-throughput data processing environments---such as network intrusion detection, financial transaction monitoring, and clinical screening---inevitably experience performance degradation due to temporal distribution shifts. In these operational settings, human-review capacity is strictly constrained by operational overhead, analyst workload, and economic cost. This study models capacity-constrained operational inspection environments using discrete 1\%, 2\%, 5\%, and 10\% review budgets. Existing selective classification frameworks prioritize audit queues relying almost exclusively on model output probabilities (confidence). However, complex classifiers frequently exhibit uncalibrated confidence under covariate and concept shift, resulting in misclassifications that maintain high posterior probabilities and bypass conventional confidence-based safety filters.

This project presents an adaptive, multi-factor human-review allocation framework designed to evaluate capacity-constrained error mitigation under distribution shift. The framework synthesizes four complementary reliability signals: (1) prediction uncertainty derived from normalized probabilities, (2) regularized Mahalanobis distance to capture geometric out-of-distribution (OOD) feature novelty, (3) Population Stability Index (PSI) to track macroscopic batch distribution drift, and (4) empirical historical error proximity mapped from validation error neighborhoods. The resulting composite risk score ranks inferences into a prioritized review queue, with a heuristic feedback mechanism adjusting signal weighting based on verified audit error density.

The proposed methodology is empirically evaluated on a 500,000-flow chronological partition from the large-scale BigFlow-NIDS benchmark (66.9 million total flows, 55 attributes), partitioned into 350,000 training, 75,000 validation, and 75,000 shifted-test instances. BigFlow-NIDS is a large-scale derived network intrusion detection benchmark constructed from publicly available NetFlow datasets. Under temporal shift, classifier error rate surges 4.55x from 1.004\% (753 errors) on validation to 4.569\% (3,427 errors) on shifted test. In offline benchmarking with a simulated ground-truth review oracle, the adaptive multi-factor allocation strategy captures 17.80\% of total shifted errors (610 errors) at a 1\% review budget (versus 13.51\% / 463 errors for pure confidence) and 34.26\% of errors (1,174 errors) at a 2\% budget (versus 26.73\% / 916 errors for pure confidence). Crucially, empirical findings reveal an honest budget-dependent crossover: while multi-factor allocation achieves higher error capture under tight review constraints (1\%--2\%), pure confidence ranking achieves higher error capture at broader budgets (5\%--10\%), capturing 65.45\% and 84.01\% of errors respectively. To date, 13 of 14 functional requirements (13 of 14 functional requirements are fully implemented; FR14 is partially implemented.) have been fully implemented across 26 sequential Python scripts, supported by an interactive local Streamlit exploratory analytics dashboard.

